% ============================================================
\section{Coherence Defect Rules (Chain and Product Laws)}
\label{sec:defect-rules}
% ============================================================

\begin{quote}\small
\textbf{Goal:} Record exact rules for how coherence defects transform under composition and bilinear operations.\\
\textbf{Key objects:} Coherence defect $\CCTau{h}{F}{x}$, bilinear defect $\CCTauBi{h}{B}{x}{y}$, leakage correction.\\
\textbf{Payoff:} These are the chain and product rules of Coherence Calculus.
\end{quote}

This section records exact algebraic rules for how coherence defects behave under composition and bilinear operations.
These are the chain and product rules of Coherence Calculus. They are
analogous to the Leibniz and chain rules of classical calculus, but stated
without derivatives or limits.
Proofs use the Resolution-of-Identity Expansion (Lemma~\ref{lem:resolution-identity}) and Block Splitting (Lemma~\ref{lem:block-splitting}) from Section~\ref{sec:foundational-lemmas}.

\subsection{Coherence defect of a map}

\begin{definition}[Coherence defect of a map]
\label{def:coherence-defect}\lean{coherenceDefect}
Let $\CCPi{h}$ be a horizon projection on a Hilbert space $\cH$ and let $\mathcal{F}:\cH \to \cH$ be a map (linear or nonlinear, defined on the relevant domain).
The \emph{coherence defect} of $\mathcal{F}$ at $x$ is
\[
\CCTau{h}{\mathcal{F}}{x} := \CCPi{h} \mathcal{F}(x) - \mathcal{F}(\CCPi{h} x).
\]
\end{definition}

\begin{remark}[Realized image of the closure defect]
\label{rem:coherence-defect-realized-closure-defect}\lean{CoherenceCalculus.coherenceDefect, CoherenceCalculus.closureDefect}
The bedrock closure defect \(\Lambda\) measures failure of Constancy to lift to
an incidence as a whole
(Definition~\ref{def:closure-defect}). Once a horizon realization has been
chosen, that abstract obstruction is seen on the realized side through the
mismatch between ``apply the law, then restrict'' and ``restrict, then apply
the realized law.'' The coherence defect \(\CCTau{h}{\mathcal F}{x}\) is
exactly that realized restriction-side image. It therefore records not the
origin of nonclosure, but its observable failure to commute with a chosen
realization map.
\end{remark}

\subsection{Scale autonomy and realized exactification}

\begin{definition}[Realized exactification generator]
\label{def:realized-exactification-generator}\lean{CoherenceCalculus.RealizedExactificationGenerator}
Let \(\cH\) carry a horizon tower \((\CCPi{h})\). A map
\[
\mathbb X:\cH\to\cH
\]
is called a \emph{realized exactification generator} if a realized selection
rule has been fixed on the chosen horizon family and there exists a
realization of incidences by states \(I\mapsto x_I\in\cH\) such that
\(\mathbb X(x_I)\) records the realized continuation selected by that rule.
Thus \(\mathbb X\) is not an arbitrary operator but a realized image of the
bedrock termination/exactification law pair relative to a chosen realized
selection of exactifying continuations on the given horizon family.
\end{definition}

\begin{definition}[Autonomous scale law]
\label{def:autonomous-scale-law}\lean{CoherenceCalculus.AutonomousScaleLaw}
Fix a realized exactification generator \(\mathbb X\). A map
\[
\mathcal L_h:\CCHobs{h}\to\CCHobs{h}
\]
is an \emph{autonomous scale law} at horizon \(h\) if it closes the realized
generator on that observed sector.
\end{definition}

\begin{definition}[Autonomy defect]
\label{def:autonomy-defect}\lean{CoherenceCalculus.autonomyDefect}
Fix a realized exactification generator \(\mathbb X\) and an autonomous scale
candidate \(\mathcal L_h\). The \emph{autonomy defect} at horizon \(h\) is
\[
\mathfrak D_h(x)
:=
\CCPi{h}\mathbb X(x)-\mathcal L_h(\CCPi{h}x).
\]
Equivalently, the scale-autonomy equation is
\[
\CCPi{h}\mathbb X = \mathcal L_h\CCPi{h},
\]
and it holds exactly precisely when \(\mathfrak D_h(x)=0\) for all \(x\) in the
domain under consideration.
\end{definition}

\begin{definition}[Autonomous horizon]
\label{def:autonomous-horizon}\lean{CoherenceCalculus.AutonomousHorizon}
A horizon \(h\) is \emph{autonomous} for the realized exactification generator
\(\mathbb X\) if there exists an autonomous scale law \(\mathcal L_h\) whose
autonomy defect vanishes identically:
\[
\mathfrak D_h(x)=0
\qquad
\text{for all }x.
\]
\end{definition}

\begin{proposition}[Autonomy criterion and uniqueness]
\label{prop:autonomy-criterion-uniqueness}\lean{CoherenceCalculus.autonomy_criterion_uniqueness, CoherenceCalculus.autonomous_scale_law_unique}
Let \(\mathbb X:\cH\to\cH\) be a realized exactification generator and fix a
horizon \(h\). Then the following are equivalent:
\begin{enumerate}[label=(\roman*), leftmargin=2em, itemsep=0.2em]
\item there exists a map \(\mathcal L_h:\CCHobs{h}\to\CCHobs{h}\) such that
      \[
      \CCPi{h}\mathbb X=\mathcal L_h\CCPi{h};
      \]
\item whenever \(\CCPi{h}x=\CCPi{h}x'\), one has
      \[
      \CCPi{h}\mathbb X(x)=\CCPi{h}\mathbb X(x').
      \]
\end{enumerate}
When these conditions hold, the autonomous scale law \(\mathcal L_h\) is
unique and is given by
\[
\mathcal L_h(y):=\CCPi{h}\mathbb X(x)
\qquad
\text{for any }x\text{ with }\CCPi{h}x=y.
\]
\end{proposition}

\begin{proof}
If \(\CCPi{h}\mathbb X=\mathcal L_h\CCPi{h}\), then equal horizon data imply
equal horizon outputs, proving \emph{(ii)}. Conversely, if \emph{(ii)} holds,
define \(\mathcal L_h\) on \(\CCHobs{h}\) by the displayed formula. The
assumption guarantees that this is well defined, and then by construction
\(\CCPi{h}\mathbb X=\mathcal L_h\CCPi{h}\). Uniqueness is immediate.
\end{proof}

\begin{definition}[Approximate autonomy]
\label{def:approximate-autonomy}\lean{CoherenceCalculus.ApproximateAutonomy}
Let \(\mathbb X\) be a realized exactification generator, let
\(\mathcal C\subseteq\cH\), and let \(\varepsilon\ge 0\). A horizon \(h\) is
\emph{\(\varepsilon\)-autonomous} on \(\mathcal C\) if there exists a map
\(\mathcal L_h:\CCHobs{h}\to\CCHobs{h}\) such that
\[
\|\CCPi{h}\mathbb X(x)-\mathcal L_h(\CCPi{h}x)\|\le \varepsilon
\qquad
(x\in\mathcal C).
\]
\end{definition}

\begin{remark}[Hydrodynamic interpretation]
\label{rem:hydrodynamic-approximate-autonomy}\lean{CoherenceCalculus.ApproximateAutonomy}
Approximate autonomy is the natural scale-law criterion for effective physics.
A hydrodynamic law exists on coarse conserved variables when the corresponding
horizon is \(\varepsilon\)-autonomous with small \(\varepsilon\); the mean free
path marks the scale below which that defect is no longer small and
hydrodynamic closure fails.
\end{remark}

\begin{proposition}[Coherence defect and autonomy defect]
\label{prop:coherence-defect-and-autonomy-defect}\lean{CoherenceCalculus.coherence_defect_and_autonomy_defect}
Let \(\mathbb X:\cH\to\cH\) be a realized exactification generator and let
\(\mathcal L_h:\CCHobs{h}\to\CCHobs{h}\) be an autonomous scale candidate at
horizon \(h\). Then for every \(x\in\cH\),
\[
\CCTau{h}{\mathbb X}{x}
\;=\;
\mathfrak D_h(x)
+
\bigl(\mathcal L_h(\CCPi{h}x)-\mathbb X(\CCPi{h}x)\bigr).
\]
In particular, if
\[
\mathcal L_h(y)=\mathbb X(y)
\qquad
\text{for every }y\in\CCHobs{h},
\]
then
\[
\CCTau{h}{\mathbb X}{x}=\mathfrak D_h(x).
\]
\end{proposition}

\begin{proof}
Add and subtract \(\mathcal L_h(\CCPi{h}x)\):
\begin{align*}
\CCTau{h}{\mathbb X}{x}
&=
\CCPi{h}\mathbb X(x)-\mathbb X(\CCPi{h}x)\\
&=
\bigl(\CCPi{h}\mathbb X(x)-\mathcal L_h(\CCPi{h}x)\bigr)
+
\bigl(\mathcal L_h(\CCPi{h}x)-\mathbb X(\CCPi{h}x)\bigr)\\
&=
\mathfrak D_h(x)
+
\bigl(\mathcal L_h(\CCPi{h}x)-\mathbb X(\CCPi{h}x)\bigr).
\end{align*}
If \(\mathcal L_h\) is the restriction of \(\mathbb X\) to \(\CCHobs{h}\), the
second term vanishes.
\end{proof}

\begin{remark}[Meaning of the scale-autonomy equation]
\label{rem:meaning-of-scale-autonomy-equation}\lean{CoherenceCalculus.AutonomousScaleLaw, CoherenceCalculus.autonomyDefect, CoherenceCalculus.AutonomousHorizon}
The equation
\[
\CCPi{h}\mathbb X = \mathcal L_h\CCPi{h}
\]
is the realized form of autonomous closure at scale \(h\). When it holds, the
observed sector has a law of its own. When it fails, unresolved structure leaks
across the horizon and the scale is still open. The coherence defect is
therefore the native operator-side witness of failure of scale autonomy.
\end{remark}

\begin{quote}\small
\textbf{Intuition.} The defect $\CCTau{h}{\mathcal{F}}{x}$ measures how much $\mathcal{F}$ fails to commute with the horizon projection.
If $\CCTau{h}{\mathcal{F}}{x} = 0$, then applying $\mathcal{F}$ before or after projecting gives the same observed result.
Nonzero defect means the observer at level $h$ cannot faithfully represent the action of $\mathcal{F}$.
\end{quote}

\paragraph{Linear vs.\ nonlinear.}
If $\mathcal{F}$ is linear and $\CCPi{h} \mathcal{F} = \mathcal{F} \CCPi{h}$, then $\CCTau{h}{\mathcal{F}}{x} = 0$ for all $x$.
For nonlinear $\mathcal{F}$, the defect is generically nonzero.
For example, if $\mathcal{F}(x) = x^2$ (pointwise square) and $\CCPi{h}$ is
block averaging, then $\CCTau{h}{\mathcal{F}}{x}$ equals the within-block
variance, which is the unresolved fluctuation contribution.

\begin{definition}[Two-argument defect for bilinear maps]
\label{def:bilinear-defect}\lean{bilinearDefect}
For a bilinear map $B:\cH \times \cH \to \cH$, define
\[
\CCTauBi{h}{B}{x}{y} := \CCPi{h} B(x,y) - B(\CCPi{h} x, \CCPi{h} y).
\]
\end{definition}

\subsection{Chain rule (composition)}

\begin{proposition}[Chain rule (exact decomposition)]
\label{prop:chain-rule}\lean{chain_rule_statement}
Let $\mathcal{F},\mathcal{G}:\cH \to \cH$ be maps and $\CCPi{h}$ a horizon projection.
Then
\[
\CCTau{h}{\mathcal{F} \circ \mathcal{G}}{x}
\;=\;
\CCTau{h}{\mathcal{F}}{\mathcal{G}(x)} \;+\; \bigl[ \mathcal{F}(\CCPi{h} \mathcal{G}(x)) - \mathcal{F}(\mathcal{G}(\CCPi{h} x)) \bigr].
\]
The first term is the defect of $\mathcal{F}$ evaluated at the full output $\mathcal{G}(x)$.
The second term is a \emph{mismatch term} measuring how $\mathcal{F}$ amplifies the defect of $\mathcal{G}$.
\end{proposition}

\begin{proof}[Proof: direct expansion]
Write
\begin{align*}
\CCTau{h}{\mathcal{F} \circ \mathcal{G}}{x}
&= \CCPi{h} \mathcal{F}(\mathcal{G}(x)) - \mathcal{F}(\mathcal{G}(\CCPi{h} x)) \\
&= \bigl[\CCPi{h} \mathcal{F}(\mathcal{G}(x)) - \mathcal{F}(\CCPi{h} \mathcal{G}(x))\bigr] + \bigl[\mathcal{F}(\CCPi{h} \mathcal{G}(x)) - \mathcal{F}(\mathcal{G}(\CCPi{h} x))\bigr] \\
&= \CCTau{h}{\mathcal{F}}{\mathcal{G}(x)} + \bigl[\mathcal{F}(\CCPi{h} \mathcal{G}(x)) - \mathcal{F}(\mathcal{G}(\CCPi{h} x))\bigr]. \qedhere
\end{align*}
\end{proof}

\paragraph{Mismatch term structure.}
The mismatch term $\mathcal{F}(\CCPi{h} \mathcal{G}(x)) - \mathcal{F}(\mathcal{G}(\CCPi{h} x))$ measures how $\mathcal{F}$ transforms the difference $\CCPi{h} \mathcal{G}(x) - \mathcal{G}(\CCPi{h} x) = \CCTau{h}{\mathcal{G}}{x}$.
If $\mathcal{F}$ is Lipschitz with constant $L$, then
\[
\|\mathcal{F}(\CCPi{h} \mathcal{G}(x)) - \mathcal{F}(\mathcal{G}(\CCPi{h} x))\| \;\le\; L \cdot \|\CCTau{h}{\mathcal{G}}{x}\|.
\]
If $\mathcal{F}$ is linear, the mismatch term equals $\mathcal{F}(\CCTau{h}{\mathcal{G}}{x})$ exactly.

\paragraph{Comparison with classical chain rule.}
In classical calculus, $(f \circ g)'(x) = f'(g(x)) \cdot g'(x)$.
Here the decomposition is additive rather than multiplicative, and the ``derivative'' is replaced by the defect functional.
The mismatch term plays the role of the $g'(x)$ factor, encoding how inner-map defects propagate through the outer map.

\subsection{Product rule (bilinear maps)}

\begin{proposition}[Product rule for bilinear $B$]
\label{prop:product-rule}\lean{product_rule_bilinear}
Let $B:\cH \times \cH \to \cH$ be bilinear and $\CCPi{h}$ a horizon projection with $\CCPibar{h} := I - \CCPi{h}$.
Assume the \emph{closure condition}: $\CCPi{h} B(\CCPi{h} x, \CCPi{h} y) = B(\CCPi{h} x, \CCPi{h} y)$, i.e., $B$ applied to two observable elements remains observable.
Then
\[
\CCTauBi{h}{B}{x}{y}
\;=\;
\CCPi{h} B(\CCPi{h} x, \CCPibar{h} y)
\;+\;
\CCPi{h} B(\CCPibar{h} x, \CCPi{h} y)
\;+\;
\CCPi{h} B(\CCPibar{h} x, \CCPibar{h} y).
\]
\end{proposition}

(Proof in Appendix~\ref{sec:appendix-tower}.)

\begin{quote}\small
\textbf{Intuition.} The bilinear defect decomposes into three ``cross terms'':
observable $\times$ shadow, shadow $\times$ observable, and shadow $\times$ shadow.
These are the contributions from unresolved components interacting.
For $B(x,y) = x \cdot y$ (pointwise product) and block averaging, this is exactly the covariance structure: the defect is the ``within-block covariance'' of $x$ and $y$.
\end{quote}

\begin{corollary}[Quadratic case]
\label{cor:quadratic-defect}\lean{quadratic_defect}
For $B(x,x)$ (e.g., $x \mapsto x^2$ or $x \mapsto \|x\|^2$),
\[
\CCTauBi{h}{B}{x}{x}
\;=\;
2\,\CCPi{h} B(\CCPi{h} x, \CCPibar{h} x)
\;+\;
\CCPi{h} B(\CCPibar{h} x, \CCPibar{h} x).
\]
If $B$ is symmetric and
$\CCPi{h} B(\CCPi{h} x, \CCPibar{h} x) = 0$ (orthogonality), then
$\CCTauBi{h}{B}{x}{x} = \CCPi{h} B(\CCPibar{h} x, \CCPibar{h} x)$, so the
defect is purely the shadow self-interaction.
\end{corollary}

\paragraph{Variance interpretation.}
For $B(x,y) = xy$ (scalar multiplication) and $\CCPi{h}$ a conditional expectation (block average), the identity $\CCTauBi{h}{B}{x}{x} = \CCPi{h}[(\CCPibar{h} x)^2]$ is the within-block variance.
This is the ``closure term'' that arises in any coarse-graining of a quadratic nonlinearity.

\subsection{Operator multiplication as a bilinear map}

A key special case of Proposition~\ref{prop:product-rule} is operator composition on $\cB(\cH)$, treated as a bilinear map.

\begin{lemma}[Operator product closure under compression]
\label{lem:operator-product-closure}\lean{leakageCorrection}
For bounded operators $A, B \in \cB(\cH)$ and horizon projection $\CCPi{h}$ with $\CCPibar{h} = I - \CCPi{h}$, define the \emph{horizon-compressed product}
\[
\CCRestr{(AB)}{h} := \CCPi{h} A B \CCPi{h}.
\]
Then
\[
\CCRestr{(AB)}{h} = \CCObs{A}{h} \CCObs{B}{h} + \CCLeakCorr{h}{A}{B},
\]
where $\CCObs{A}{h} = \CCRestr{A}{h}$, $\CCObs{B}{h} = \CCRestr{B}{h}$, and the \emph{leakage correction} is
\[
\CCLeakCorr{h}{A}{B} := \CCPi{h} A \CCPibar{h} B \CCPi{h}.
\]
\end{lemma}

\begin{proof}
Insert $I = \CCPi{h} + \CCPibar{h}$ between $A$ and $B$:
\[
\CCPi{h} A B \CCPi{h} = \CCPi{h} A (\CCPi{h} + \CCPibar{h}) B \CCPi{h} = \CCPi{h} A \CCPi{h} B \CCPi{h} + \CCPi{h} A \CCPibar{h} B \CCPi{h} = \CCObs{A}{h} \CCObs{B}{h} + \CCLeakCorr{h}{A}{B}. \qedhere
\]
\end{proof}

\begin{quote}\small
\textbf{Intuition.} The compressed product $\CCRestr{(AB)}{h}$ is not simply the product of compressed operators $\CCObs{A}{h} \CCObs{B}{h}$. The correction $\CCLeakCorr{h}{A}{B}$ captures ``round trips through the shadow'': the intermediate result $B \CCPi{h} x$ may have components in the shadow ($\CCPibar{h} B \CCPi{h} x$), which $A$ then maps back into the observable.
\end{quote}

\begin{corollary}[Connection to Proposition~\ref{prop:product-rule}]
\label{cor:product-connection}\lean{operator_mult_defect_eq_leakage}
Define the bilinear map $M: \cB(\cH) \times \cB(\cH) \to \cB(\cH)$ by $M(A, B) = AB$. The coherence defect for operator multiplication satisfies
\[
\CCTauBi{h}{M}{A}{B} = \CCPi{h} (AB) - (\CCRestr{A}{h})(\CCRestr{B}{h}) = \CCLeakCorr{h}{A}{B}.
\]
This is exactly the correction term from Lemma~\ref{lem:operator-product-closure}.
\end{corollary}

\paragraph{Relation to defect operator.}
When $A = B = d$ is a nilpotent differential, the leakage correction $\CCLeakCorr{h}{d}{d} = \CCDefectExpr{d}{h} = \CCDefect{d}{h}$ is precisely the defect operator from the Horizon Fundamental Theorem (Section~\ref{sec:hft}).

\subsection{Defect energy under bilinear maps}

\begin{lemma}[Defect energy bound for bilinear maps]
\label{lem:bilinear-energy}\lean{bilinear_defect_bound_statement}
Let $B:\cH \times \cH \to \cH$ be a bounded bilinear map with $\|B(x,y)\| \le M\|x\|\,\|y\|$.
Then
\[
\|\CCTauBi{h}{B}{x}{y}\| \;\le\; M \bigl( \|\CCPi{h} x\|\,\|\CCPibar{h} y\| + \|\CCPibar{h} x\|\,\|\CCPi{h} y\| + \|\CCPibar{h} x\|\,\|\CCPibar{h} y\| \bigr).
\]
In particular, if $\|\CCPibar{h} x\| \le \epsilon$ and $\|\CCPibar{h} y\| \le \epsilon$, then $\|\CCTauBi{h}{B}{x}{y}\| = O(\epsilon)$.
\end{lemma}

\begin{proof}
Apply the triangle inequality to Proposition~\ref{prop:product-rule} and use $\|\CCPi{h}\| \le 1$.
\end{proof}

\subsection{Kernel defect under nested horizons}

The following lemma describes how kernel defects behave under horizon refinement.

\begin{lemma}[Kernel monotonicity under domain restriction]
\label{lem:kernel-monotonicity}\lean{kernel_monotone}
Let $(\CCPi{h})$ be a transitive horizon tower with $\CCPi{h} \CCPi{h'} = \CCPi{h}$ for $h \le h'$.
Let $C:\cH \to \cH'$ be a constraint map.
Define the \emph{domain-restricted constraint map} $C_h := C|_{\CCHobs{h}}$.
\begin{enumerate}[label=(\roman*)]
\item $\Ker(C_h) = \Ker(C) \cap \CCHobs{h}$.
\item For nested horizons $h \le h'$: $\Ker(C_h) \subseteq \Ker(C_{h'})$, and hence $\dim\Ker(C_h) \le \dim\Ker(C_{h'})$.
\end{enumerate}
\end{lemma}

\begin{proof}
(i) If $x \in \CCHobs{h}$, then $C_h x = C x$, so $C_h x = 0$ iff $C x = 0$.

(ii) If $x \in \Ker(C_h) = \Ker(C) \cap \CCHobs{h}$, then $x \in \Ker(C)$ and $x \in \CCHobs{h} \subseteq \CCHobs{h'}$, so $x \in \Ker(C) \cap \CCHobs{h'} = \Ker(C_{h'})$.
\end{proof}

\paragraph{Monotonicity direction.}
The direction of kernel monotonicity depends on whether horizon restriction is applied to the domain or codomain of $C$.
For domain restriction (as above), finer horizons have larger kernels.
For codomain restriction ($C_h = R_h \circ C$ where $R_h$ restricts the target), the direction may reverse.
We record the domain-restriction case as the default.

\subsection{Calculator lemma: carrying resolution terms for linear operators}
\label{subsec:calculator-lemma}

This lemma provides a practical recipe for computing with defects when applying linear operators under finite observability.

\begin{lemma}[Resolution correction for linear operators]
\label{lem:resolution-correction}\lean{resolution_correction}
Let $L: \cH \to \cH$ be a bounded linear operator (or closed with appropriate domain) and $\CCPi{h}$ a horizon projection. For any $x \in \cH$,
\[
\CCPi{h} L x = (\CCRestr{L}{h}) x + \CCPi{h} L \CCPibar{h} x.
\]
Equivalently:
\[
\text{(projected truth)} = \text{(observed computation)} + \text{(shadow correction)}.
\]
\end{lemma}

\begin{proof}
Insert $I = \CCPi{h} + \CCPibar{h}$ before $x$:
\[
\CCPi{h} L x = \CCPi{h} L (\CCPi{h} + \CCPibar{h}) x = \CCPi{h} L \CCPi{h} x + \CCPi{h} L \CCPibar{h} x. \qedhere
\]
\end{proof}

\paragraph{Computational interpretation.}
The observed operator $\CCObs{L}{h} = \CCRestr{L}{h}$ acting on the observed state $\CCPi{h} x$ gives only part of the projected truth. The shadow correction $\CCPi{h} L \CCPibar{h} x$ accounts for how $L$ maps the unobserved part of $x$ into the observable subspace.

\begin{lemma}[Finite-rank computation recipe]
\label{lem:finite-rank-recipe}\lean{finite_rank_observable}
Suppose $\CCPi{h}$ has finite rank with orthonormal basis $\{\psi_1, \ldots, \psi_r\}$ for $\CCHobs{h}$.
Define coefficient vectors $c(x) \in \mathbb{R}^r$ by $c_i(x) := \langle x, \psi_i \rangle$.
Then:
\begin{enumerate}[label=(\roman*)]
\item The horizon-restricted operator $\CCObs{L}{h} = \CCRestr{L}{h}$ has matrix representation $(\CCObs{L}{h})_{ij} = \langle \psi_i, L \psi_j \rangle$ (the Gram matrix of $L$ in the basis).
\item The projected result is
\[
c(\CCPi{h} L x) = \CCObs{L}{h} \cdot c(\CCPi{h} x) + \delta_h(x),
\]
where the defect vector $\delta_h(x)_i = \langle \psi_i, L \CCPibar{h} x \rangle$ encodes the shadow correction in coordinates.
\item If $x \in \CCHobs{h}$ (i.e., $\CCPibar{h} x = 0$), then $\delta_h(x) = 0$ and $c(\CCPi{h} L x) = \CCObs{L}{h} \cdot c(x)$.
\end{enumerate}
\end{lemma}

\begin{proof}
(i) For $y = \sum_j c_j \psi_j \in \CCHobs{h}$, $(\CCObs{L}{h} y)_i = \langle \psi_i, \CCObs{L}{h} y \rangle = \langle \psi_i, \CCRestr{L}{h} y \rangle = \langle \psi_i, L y \rangle = \sum_j c_j \langle \psi_i, L \psi_j \rangle$.

(ii) Apply Lemma~\ref{lem:resolution-correction} and project onto the basis.

(iii) Immediate from (ii) with $\CCPibar{h} x = 0$.
\end{proof}

\begin{quote}\small
\textbf{By-hand usability.} For finite-rank horizons, represent $\CCObs{L}{h}$ as an $r \times r$ matrix by computing inner products $\langle \psi_i, L \psi_j \rangle$. The defect vector $\delta_h(x)$ measures how much information from the shadow ``leaks back in'' when $L$ is applied. This recipe reduces infinite-dimensional computations to linear algebra plus controlled error terms.
\end{quote}

\paragraph{Further developments.}
The defect rules in this section provide the basic calculus of coherence defects.
Subsequent sections extend this machinery:
\begin{itemize}[itemsep=0.2em]
\item Section~\ref{sec:constraint-compiler} uses these rules within the admissible-completion selection scheme.
\item Section~\ref{sec:resolvent-defect} provides exact resummation via the Feshbach--Schur identity.
\item Section~\ref{sec:coherence-gap} derives spectral bounds from defect positivity.
\item Section~\ref{sec:log-periodic} specializes to log-scale horizon families.
\item Section~\ref{sec:advanced-operators} packages scale flow, spectral selection, trace aggregation, and canonical leakage direction.
\end{itemize}
